\documentclass[12pt,a4paper]{article}
\usepackage[utf8]{inputenc}
\usepackage[spanish,es-noquoting]{babel}
\usepackage[T1]{fontenc}
\usepackage{lmodern}
\usepackage{geometry}
\geometry{top=2cm, bottom=2cm, left=2.5cm, right=2.5cm}
\usepackage{xcolor}
\usepackage{listings}
\usepackage{enumitem}
\usepackage{fancyhdr}
\usepackage{hyperref}

\definecolor{codegreen}{rgb}{0,0.6,0}
\definecolor{backcolour}{rgb}{0.95,0.95,0.92}

\lstset{
    language=C,
    backgroundcolor=\color{backcolour},
    commentstyle=\color{codegreen},
    keywordstyle=\color{blue},
    basicstyle=\ttfamily\small,
    breaklines=true,
    frame=single,
    numbers=left,
    tabsize=4,
    extendedchars=true,
    showstringspaces=false,
    literate={á}{{\'a}}1 {é}{{\'e}}1 {í}{{\'i}}1 {ó}{{\'o}}1 {ú}{{\'u}}1 {ñ}{{\~n}}1
}

\pagestyle{fancy}
\fancyhf{}
\lhead{Monitorías Sistemas Operativos 2026-I}
\rhead{Capítulo 0: Archivos}
\cfoot{\thepage}

\title{\textbf{Guía de Aprendizaje: Manejo de Archivos en C}}
\author{Malak Sanchez \\ \small Monitorías de Sistemas Operativos}
\date{}

\begin{document}

\maketitle

\section{Motivación}
Los datos en memoria RAM son volátiles: desaparecen al apagar el equipo. En Sistemas Operativos, necesitamos persistencia para guardar configuraciones, logs o bases de datos. Para ello usamos archivos.

\section{Idea Fundamental}
Un archivo es un flujo (\textit{stream}) de datos almacenados en un dispositivo de almacenamiento. C maneja archivos mediante una estructura especial llamada \texttt{FILE}.

\section{Sintaxis Básica}
\begin{lstlisting}[numbers=none]
FILE *file_ptr = fopen("filename.txt", "mode");
fclose(file_ptr);
\end{lstlisting}

\section{Modos de Apertura}
El segundo argumento de \texttt{fopen} define qué podemos hacer con el archivo. Aquí los más comunes:

\begin{center}
\renewcommand{\arraystretch}{1.5}
\begin{tabular}{|c|l|c|c|}
\hline
\textbf{Modo} & \textbf{Descripción} & \textbf{¿Crea archivo?} & \textbf{¿Sobrescribe?} \\ \hline
\texttt{``r''} & \textbf{Lectura:} El archivo debe existir. & No & No \\ \hline
\texttt{``w''} & \textbf{Escritura:} Crea el archivo o lo vacía si existe. & \textbf{Sí} & \textbf{Sí} \\ \hline
\texttt{``a''} & \textbf{Anexar:} Escribe al final del archivo. & \textbf{Sí} & No \\ \hline
\texttt{``r+''} & \textbf{Lectura/Escritura:} El archivo debe existir. & No & No \\ \hline
\texttt{``w+''} & \textbf{Lectura/Escritura:} Crea o sobrescribe. & \textbf{Sí} & \textbf{Sí} \\ \hline
\texttt{``a+''} & \textbf{Lectura/Anexar:} Lee y escribe al final. & \textbf{Sí} & No \\ \hline
\end{tabular}
\end{center}

\begin{itemize}
    \item \textbf{Modo Binario:} Si trabajas con archivos que no son de texto (imágenes, ejecutables), añade una \texttt{``b''} al modo (ej: \texttt{``rb''}, \texttt{``wb''}).
\end{itemize}

\section{Ejemplo Mínimo Funcional}
\begin{lstlisting}
#include <stdio.h>

int main() {
    // Abrir para escribir (si existe, se borra el contenido anterior)
    FILE *f = fopen("output.txt", "w");
    if(f == NULL){
        perror("Error opening file");
        return 1;
    }

    fprintf(f, "Operating Systems 2026-I\n");
    
    fclose(f);
    return 0;
}
\end{lstlisting}

\section{Explicación Paso a Paso}
\begin{enumerate}
    \item Se declara un puntero de tipo \texttt{FILE}.
    \item \texttt{fopen} abre el archivo. El modo \texttt{"w"} es para escritura (\textit{write}).
    \item Se verifica que el archivo se haya abierto correctamente (\texttt{f != NULL}).
    \item \texttt{fprintf} escribe datos en el archivo, similar a \texttt{printf}.
    \item \texttt{fclose} cierra el flujo y asegura que los datos se guarden en el disco.
\end{enumerate}

\section{Qué ocurre en memoria}
El Sistema Operativo mantiene una tabla de descriptores de archivos para cada proceso. El puntero \texttt{FILE} apunta a una estructura que contiene el estado del flujo (posición actual, buffer, etc.).

\section{Patrones comunes de uso}
Lectura línea por línea:
\begin{lstlisting}
char buffer[100];
while(fgets(buffer, 100, f) != NULL){
    printf("%s", buffer);
}
\end{lstlisting}

\section{Errores Comunes}
\begin{itemize}
    \item \textbf{No cerrar el archivo:} Puede causar pérdida de datos o agotar el límite de archivos abiertos del sistema operativo.
    \item \textbf{Modos incorrectos:} Intentar escribir en un archivo abierto como \texttt{``r''} (lectura).
\end{itemize}

\section{Buenas Prácticas}
\begin{itemize}
    \item Siempre verificar si \texttt{fopen} tuvo éxito (usar \texttt{perror} ayuda a debuguear).
    \item Cerrar los archivos apenas se termine de usarlos.
\end{itemize}

\section{Ejercicios}
\begin{enumerate} 
    \item Cree un programa que lea un archivo llamado \texttt{data.txt} e imprima cuántas líneas tiene.
    \item Escriba un programa que pida el nombre del usuario y lo guarde en un archivo \texttt{config.log}.
\end{enumerate}

\end{document}
